\section{StudyBirds}
\label{sec:system}

The system is organized around a course, its learning topics, and
persistent study sessions. The description below combines the student
prototype with the intended instructor feedback workflow; validation
of the complete loop is part of the evaluation plan.

\subsection{Course Materials and Study Context}

Course preparation organizes uploaded materials into a course map and
learning units with links to their sources. The intended review process
lets an instructor inspect the organization and proposed learning
activities before publishing them. Students may also create private
courses for their own study. Such material and its associated learning
history should remain separate from instructor-visible class reporting.

\subsection{Guided Study}

Students enter a topic with a purpose such as learning, practice, review,
or exam preparation. A conversational tutor works alongside a board
containing explanations, questions, and interactive activities.
The tutor can respond to an attempt, provide a hint, or revisit a
prerequisite. Generated examples should be distinguished from original
course content so that students can check the source when an explanation
is unclear.

The student prototype supports pausing and resuming a session, retaining
its teaching content, and producing a session recap and a topic study
sheet. These outputs are intended to help a learner decide where to
continue. They must also indicate which ideas were discussed and which
were actually checked.

\subsection{Evidence and Assistance}

An answer alone does not describe the conditions under which it was
produced. StudyBirds records attempts together with assistance context
and uses that distinction when presenting learning evidence. For example,
solving a question after an explicit hint may justify a follow-up
independent check, but should not be reported as independent success.
Session summaries should preserve this distinction and retain the
history of incorrect attempts and subsequent corrections.

This mechanism is intended to make progress claims inspectable.
Its validity depends on correct grading, accurate identification of
assistance, and a defensible connection between an activity and the
learning objective it assesses.

\subsection{Instructor Feedback}

The instructor-facing goal is to turn permitted class evidence into a
concise explanation of a specific difficulty, with a link to the relevant
course source. For example, an insight might distinguish confusion about
a formula's assumptions from an arithmetic error when applying it.
An instructor could then inspect the evidence and revise an explanation
or prepare a targeted activity. The usefulness and privacy of this
workflow require separate validation; removing names alone does not
establish anonymity.
